\section{Banco de questões — Legislação Específica do TCE-MA}

\subsection*{N1 — Fundamentos em nível de concurso}

\begin{questao}{\nivel{N1} 07.01 — Cebraspe (C/E)}
O fato de o TCE-MA auxiliar a Assembleia Legislativa no controle externo significa que o Tribunal se encontra hierarquicamente subordinado ao Poder Legislativo estadual.
\end{questao}

\begin{questao}{\nivel{N1} 07.02 — Cebraspe (C/E)}
A Lei Orgânica do TCE-MA e o Regimento Interno ocupam posições normativas distintas, de modo que regra regimental não pode validamente afastar comando estabelecido em lei.
\end{questao}

\begin{questao}{\nivel{N1} 07.03 — FGV}
No sistema de controle externo, a apreciação das contas anuais do chefe do Poder Executivo distingue-se do julgamento de contas de responsáveis porque, em regra, a primeira resulta em parecer prévio, ao passo que o segundo produz deliberação de julgamento pelo Tribunal de Contas. Assinale a alternativa que melhor traduz essa distinção.
\begin{enumerate}[label=\Alph*)]
\item Todo parecer prévio equivale a condenação definitiva do gestor.
\item O parecer prévio integra processo de apreciação das contas de governo; o julgamento alcança contas de responsáveis submetidos à jurisdição de contas.
\item O Tribunal de Contas apenas emite recomendações sem efeitos jurídicos.
\item As contas de governo e as contas de gestão são categorias necessariamente idênticas.
\item O parecer prévio substitui sempre a deliberação do Poder Legislativo.
\end{enumerate}
\end{questao}

\begin{questao}{\nivel{N1} 07.04 — Cebraspe (C/E)}
Unidade técnica de fiscalização pode produzir instrução e proposta de encaminhamento, mas sua manifestação não se confunde com a deliberação do órgão competente do Tribunal.
\end{questao}

\begin{questao}{\nivel{N1} 07.05 — FCC}
Em processo de controle externo, a existência de dano ao erário e a aplicação de multa são institutos distintos. É correto afirmar que:
\begin{enumerate}[label=\Alph*)]
\item débito e multa são sinônimos.
\item débito relaciona-se à recomposição de dano quantificado, enquanto multa possui natureza sancionatória.
\item multa somente pode ser aplicada quando houver débito.
\item débito dispensa identificação do responsável.
\item multa substitui automaticamente o ressarcimento.
\end{enumerate}
\end{questao}

\begin{questao}{\nivel{N1} 07.06 — Cebraspe (C/E)}
A tomada de contas especial tem por finalidade, entre outros aspectos, apurar fatos, identificar responsáveis e quantificar dano quando presentes os pressupostos de sua instauração.
\end{questao}

\begin{questao}{\nivel{N1} 07.07 — Cebraspe (C/E)}
A Instrução Normativa TCE-MA nº 82/2025 disciplina fiscalização, acompanhamento e julgamento da execução de emendas parlamentares estaduais e municipais, com ênfase em transparência, rastreabilidade e conformidade constitucional.
\end{questao}

\begin{questao}{\nivel{N1} 07.08 — Vunesp}
A Lei Estadual nº 9.936/2013, indicada no edital, deve ser estudada principalmente como norma relativa:
\begin{enumerate}[label=\Alph*)]
\item ao Código Tributário estadual.
\item à organização administrativa do TCE-MA.
\item ao processo penal militar.
\item às eleições municipais.
\item à organização do Ministério Público da União.
\end{enumerate}
\end{questao}

\begin{questao}{\nivel{N1} 07.09 — Cebraspe (C/E)}
Cautelar e sanção possuem funções distintas: a cautelar visa preservar a utilidade da atuação de controle ou evitar dano, enquanto a sanção pressupõe responsabilização nos termos legais.
\end{questao}

\begin{questao}{\nivel{N1} 07.10 — Cebraspe (C/E)}
O Ministério Público de Contas integra a instrução e o controle externo segundo suas atribuições próprias, mas não se confunde com unidade técnica de auditoria do Tribunal.
\end{questao}

\subsection*{N2 — Aplicação e distinções}

\begin{questao}{\nivel{N2} 07.11 — Cebraspe (C/E)}
Se relatório técnico conclui pela irregularidade de determinada despesa, essa conclusão, por si só, equivale ao julgamento definitivo do responsável pelo TCE-MA.
\end{questao}

\begin{questao}{\nivel{N2} 07.12 — FGV}
Em processo no TCE-MA, identificou-se possível dano, mas o responsável ainda não foi chamado para se manifestar. Antes de decisão definitiva gravosa, a providência compatível com devido processo é:
\begin{enumerate}[label=\Alph*)]
\item dispensar qualquer comunicação porque Tribunais de Contas não observam contraditório.
\item assegurar oportunidade de defesa, sem prejuízo de cautelar urgente quando juridicamente cabível.
\item aplicar imediatamente débito e multa e somente depois instaurar processo.
\item remeter obrigatoriamente o caso ao Judiciário e arquivar o processo de contas.
\item substituir a defesa por parecer da unidade técnica.
\end{enumerate}
\end{questao}

\begin{questao}{\nivel{N2} 07.13 — Cebraspe (C/E)}
A competência do relator e a competência do colegiado podem coexistir no mesmo processo, cabendo ao Regimento Interno definir quais atos são monocráticos e quais dependem de deliberação colegiada.
\end{questao}

\begin{questao}{\nivel{N2} 07.14 — FCC}
Ao estudar recurso previsto na legislação do TCE-MA, o método mais seguro é relacionar simultaneamente:
\begin{enumerate}[label=\Alph*)]
\item nome do recurso, legitimidade, decisão recorrível, prazo, efeito e órgão competente.
\item apenas o prazo, pois os demais elementos são invariáveis.
\item somente a denominação do recurso.
\item regras do TCU, aplicadas automaticamente ao TCE-MA.
\item regras do CPC, independentemente do regime próprio.
\end{enumerate}
\end{questao}

\begin{questao}{\nivel{N2} 07.15 — Cebraspe (C/E)}
A existência de impropriedade formal em prestação de contas não implica, por definição e sem exame de materialidade e gravidade, que as contas devam ser classificadas como irregulares.
\end{questao}

\begin{questao}{\nivel{N2} 07.16 — Cebraspe (C/E)}
Na tomada de contas especial, a quantificação do dano é relevante porque permite vincular a irregularidade a possível obrigação de ressarcimento, desde que demonstrados responsabilidade e nexo causal.
\end{questao}

\begin{questao}{\nivel{N2} 07.17 — FGV}
Município não comprova a correta aplicação de recursos repassados e não saneia as irregularidades após medidas administrativas. À luz da lógica da IN TCE-MA nº 50/2017, a tomada de contas especial deve ser entendida como instrumento destinado a:
\begin{enumerate}[label=\Alph*)]
\item substituir o orçamento anual.
\item apurar fatos, identificar responsáveis, quantificar dano e buscar ressarcimento.
\item conceder anistia automática ao gestor.
\item julgar ação penal.
\item aprovar lei municipal.
\end{enumerate}
\end{questao}

\begin{questao}{\nivel{N2} 07.18 — Cebraspe (C/E)}
Segundo a disciplina operacional associada à IN TCE-MA nº 50/2017, a autoridade administrativa competente possui papel na instauração da tomada de contas especial, não sendo a instauração ato reservado exclusivamente ao Plenário do TCE-MA.
\end{questao}

\begin{questao}{\nivel{N2} 07.19 — Cebraspe (C/E)}
A IN nº 82/2025 foi editada em contexto de reforço da transparência e rastreabilidade de emendas parlamentares, inclusive diante das discussões constitucionais associadas à ADPF 854.
\end{questao}

\begin{questao}{\nivel{N2} 07.20 — FGV}
Na execução de emenda parlamentar, o sistema registra apenas o valor transferido e o ente beneficiário, sem indicar parlamentar proponente, objeto, plano de trabalho, cronograma ou documentos da despesa. Sob a ótica da IN nº 82/2025, o principal problema é:
\begin{enumerate}[label=\Alph*)]
\item excesso de transparência.
\item insuficiência de rastreabilidade da origem, finalidade e execução do recurso.
\item falta de competência tributária do município.
\item ausência de licença de software proprietário.
\item impossibilidade de qualquer fiscalização por Tribunal de Contas.
\end{enumerate}
\end{questao}

\begin{questao}{\nivel{N2} 07.21 — Cebraspe (C/E)}
A transparência exigida para emendas parlamentares deve permitir acompanhar o recurso ao longo do ciclo de execução, e não apenas divulgar o valor global autorizado.
\end{questao}

\begin{questao}{\nivel{N2} 07.22 — FCC}
Se norma regimental estabelece procedimento incompatível com posterior alteração da Lei Orgânica, o intérprete deve:
\begin{enumerate}[label=\Alph*)]
\item aplicar o Regimento por ser norma mais detalhada, ainda que contrarie lei.
\item harmonizar o procedimento com a lei superior e reconhecer a prevalência da disciplina legal.
\item ignorar a Lei Orgânica.
\item aplicar apenas costume administrativo.
\item transferir a competência ao Poder Judiciário automaticamente.
\end{enumerate}
\end{questao}

\begin{questao}{\nivel{N2} 07.23 — Cebraspe (C/E)}
A modificação do Regimento Interno depende do procedimento previsto no próprio Regimento e não pode ser realizada informalmente por unidade administrativa.
\end{questao}

\begin{questao}{\nivel{N2} 07.24 — Cebraspe (C/E)}
Parecer prévio sobre contas de governo e decisão de julgamento de contas de gestão possuem a mesma natureza jurídica e são sempre proferidos pelo mesmo destinatário institucional.
\end{questao}

\begin{questao}{\nivel{N2} 07.25 — FGV}
Em fiscalização de emendas parlamentares, qual evidência melhor atende à ideia de rastreabilidade?
\begin{enumerate}[label=\Alph*)]
\item planilha sem identificação da fonte do recurso.
\item cadeia que conecte emenda, proponente, plano de trabalho, empenho, liquidação, pagamento, documento fiscal e entrega física.
\item declaração genérica de regularidade.
\item relatório sem código ou processo de origem.
\item apenas fotografia da obra, sem vínculo documental.
\end{enumerate}
\end{questao}

\subsection*{N3 — Nível de prova e integração normativa}

\begin{questao}{\nivel{N3} 07.26 — Cebraspe (C/E)}
Em razão de sua autonomia funcional, o TCE-MA pode criar por regimento hipótese de sanção pecuniária sem previsão legal, desde que aprovada pela maioria absoluta de seus conselheiros.
\end{questao}

\begin{questao}{\nivel{N3} 07.27 — FGV}
Considere que uma equipe técnica identifique dano de R\$ 800 mil, mas não demonstre quem praticou a conduta nem o nexo entre comportamento e prejuízo. Para imputação de débito, o processo apresenta deficiência porque:
\begin{enumerate}[label=\Alph*)]
\item a quantificação do dano é irrelevante.
\item responsabilidade patrimonial exige suporte probatório quanto à conduta, ao responsável e ao nexo causal, além do dano.
\item débito é sanção objetiva e independe de qualquer vínculo subjetivo.
\item toda perda financeira deve ser atribuída ao chefe do Poder Executivo.
\item a unidade técnica pode dispensar instrução quando o valor for elevado.
\end{enumerate}
\end{questao}

\begin{questao}{\nivel{N3} 07.28 — Cebraspe (C/E)}
A adoção de cautelar antes da oitiva do responsável pode ser compatível com o devido processo quando houver fundamento para urgência e posterior oportunidade de contraditório, não se confundindo com julgamento definitivo.
\end{questao}

\begin{questao}{\nivel{N3} 07.29 — Cebraspe (C/E)}
No estudo de legislação específica, a utilização de prazo extraído de versão antiga do Regimento pode produzir erro mesmo quando o instituto processual permaneceu com a mesma denominação.
\end{questao}

\begin{questao}{\nivel{N3} 07.30 — FCC}
Um servidor precisa decidir qual norma consultar para saber se determinada matéria é de competência do Plenário, de Câmara ou do relator. A fonte imediatamente mais adequada, sem prejuízo da Lei Orgânica, é:
\begin{enumerate}[label=\Alph*)]
\item Regimento Interno vigente.
\item Código Penal.
\item Lei de Diretrizes Orçamentárias da União.
\item regulamento de licitações de empresa privada.
\item norma técnica ISO 9001.
\end{enumerate}
\end{questao}

\begin{questao}{\nivel{N3} 07.31 — Cebraspe (C/E)}
A publicidade de informações sobre emendas parlamentares pode ser considerada suficiente mesmo que impossibilite relacionar o valor pago ao objeto, ao beneficiário, ao plano de trabalho e aos documentos de execução, desde que o total anual seja divulgado.
\end{questao}

\begin{questao}{\nivel{N3} 07.32 — FGV}
Uma emenda transfere R\$ 2 milhões para obra municipal. O portal divulga proponente, valor e objeto, mas omite medições, notas fiscais, pagamentos e execução física. A conclusão mais adequada é:
\begin{enumerate}[label=\Alph*)]
\item a rastreabilidade está completa, pois basta conhecer o autor da emenda.
\item há lacuna de transparência da fase executória, impedindo vincular recurso financeiro a entrega e despesa comprovada.
\item a omissão é irrelevante para controle externo.
\item a transparência somente começa após o encerramento da obra.
\item documentos da execução nunca podem ser publicados.
\end{enumerate}
\end{questao}

\begin{questao}{\nivel{N3} 07.33 — Cebraspe (C/E)}
A IN nº 82/2025 alcança o ciclo da política pública financiada por emenda, envolvendo planejamento, formalização de planos de trabalho, execução, prestação de informações, transparência e controle.
\end{questao}

\begin{questao}{\nivel{N3} 07.34 — Cebraspe (C/E)}
A organização administrativa prevista na Lei nº 9.936/2013 deve ser estudada separadamente da distribuição processual do Regimento, embora os dois diplomas possam se relacionar na compreensão da estrutura do Tribunal.
\end{questao}

\begin{questao}{\nivel{N3} 07.35 — FGV}
Ao receber achado de auditoria indicando pagamento sem comprovação do objeto, o relator determina diligências para obtenção de documentos antes de submeter o mérito ao colegiado. Esse comportamento evidencia que:
\begin{enumerate}[label=\Alph*)]
\item instrução e julgamento são necessariamente o mesmo ato.
\item a formação de evidência pode anteceder a decisão de mérito e ser conduzida por atos processuais do relator.
\item o colegiado perdeu competência para decidir.
\item o relatório técnico substituiu o contraditório.
\item toda diligência equivale a sanção.
\end{enumerate}
\end{questao}

\begin{questao}{\nivel{N3} 07.36 — Cebraspe (C/E)}
Em processo com possível débito, a defesa deve poder enfrentar os fatos, a quantificação do dano e a imputação de responsabilidade, pois contraditório meramente formal não satisfaz a lógica do devido processo.
\end{questao}

\begin{questao}{\nivel{N3} 07.37 — Cebraspe (C/E)}
A existência de recomendação anterior do TCE-MA torna automaticamente desnecessária qualquer nova instrução sobre fatos posteriores, ainda que haja indícios de dano em exercício diverso.
\end{questao}

\begin{questao}{\nivel{N3} 07.38 — FCC}
Em tomada de contas especial, qual sequência é mais coerente com a finalidade do instrumento?
\begin{enumerate}[label=\Alph*)]
\item fato \textrightarrow responsável \textrightarrow dano \textrightarrow nexo \textrightarrow ressarcimento/decisão.
\item sanção \textrightarrow fato \textrightarrow defesa.
\item recurso \textrightarrow orçamento \textrightarrow julgamento penal.
\item parecer prévio \textrightarrow licitação \textrightarrow eleição.
\item recomendação \textrightarrow prescrição automática \textrightarrow arquivamento.
\end{enumerate}
\end{questao}

\begin{questao}{\nivel{N3} 07.39 — Cebraspe (C/E)}
A IN nº 50/2017 disciplina tomada de contas especial, enquanto a IN nº 82/2025 possui objeto ligado às emendas parlamentares; confundir o objeto dos dois atos é erro material de legislação específica.
\end{questao}

\begin{questao}{\nivel{N3} 07.40 — FGV}
Uma questão afirma que “todo ato de fiscalização do TCE-MA é necessariamente julgamento de contas”. A afirmação é incorreta porque:
\begin{enumerate}[label=\Alph*)]
\item fiscalização pode assumir formas e produzir providências distintas, como auditoria, inspeção, acompanhamento, determinação, recomendação ou instrução de processo.
\item o TCE-MA não fiscaliza.
\item somente o Judiciário pode examinar despesa pública.
\item contas públicas não estão sujeitas a controle externo.
\item toda fiscalização gera exclusivamente parecer prévio.
\end{enumerate}
\end{questao}

\subsection*{N4 — Avançado, controle e auditoria}

\begin{questao}{\nivel{N4} 07.41 — Cebraspe (C/E)}
Se um sistema de transparência de emendas permite alterar registros históricos sem versionamento ou trilha de auditoria, a mera presença de todos os campos exigidos é suficiente para assegurar rastreabilidade material.
\end{questao}

\begin{questao}{\nivel{N4} 07.42 — FGV}
Em auditoria sobre emendas PIX, o órgão comprova transferência e plano de trabalho, mas não consegue demonstrar correspondência entre pagamentos, fornecedores, documentos fiscais e metas entregues. O risco central é:
\begin{enumerate}[label=\Alph*)]
\item impossibilidade de verificar a cadeia de execução e o nexo entre recurso e resultado.
\item excesso de segregação de funções.
\item violação do princípio da anualidade eleitoral apenas.
\item ausência de competência do TCE para examinar recursos públicos.
\item redundância excessiva de informação.
\end{enumerate}
\end{questao}

\begin{questao}{\nivel{N4} 07.43 — Cebraspe (C/E)}
Em controle externo orientado por evidências, transparência e rastreabilidade são conceitos complementares: publicar dados não basta se eles não permitirem reconstruir a origem, o fluxo e a aplicação do recurso.
\end{questao}

\begin{questao}{\nivel{N4} 07.44 — FGV}
Um gestor sustenta que a tomada de contas especial deve ser instaurada sempre que houver qualquer erro formal, ainda que não exista dano nem necessidade de identificar responsável. A melhor análise é:
\begin{enumerate}[label=\Alph*)]
\item correta, porque a TCE é procedimento universal para toda impropriedade.
\item incorreta; a TCE possui finalidade específica ligada à apuração de fatos, responsáveis e dano, não sendo substituto indiscriminado de todo mecanismo de controle.
\item correta apenas se o erro for de digitação.
\item correta porque contraditório é dispensável.
\item incorreta porque TCE somente existe no âmbito federal.
\end{enumerate}
\end{questao}

\begin{questao}{\nivel{N4} 07.45 — Cebraspe (C/E)}
A responsabilização em processo de contas deve distinguir falha de controle, irregularidade, dano e conduta individual, evitando inferir automaticamente responsabilidade pessoal apenas da posição hierárquica ocupada.
\end{questao}

\begin{questao}{\nivel{N4} 07.46 — FGV}
O banco de dados de emendas contém código único, proponente, objeto, valor, beneficiário e cronograma, mas os registros de pagamento podem ser excluídos por administrador sem log imutável. Para auditoria, a principal recomendação é:
\begin{enumerate}[label=\Alph*)]
\item eliminar os códigos únicos.
\item implementar trilha de auditoria, controle de privilégios, versionamento e preservação dos registros de alteração/exclusão.
\item permitir exclusão anônima para agilizar correções.
\item substituir dados estruturados por documentos sem índice.
\item restringir toda informação ao ambiente interno, inclusive a que deve ser transparente.
\end{enumerate}
\end{questao}

\begin{questao}{\nivel{N4} 07.47 — Cebraspe (C/E)}
A análise de conformidade com a IN nº 82/2025 pode exigir cruzamento entre dados orçamentários, financeiros, instrumentos jurídicos, documentos fiscais e execução física, e não apenas leitura isolada do ato de emenda.
\end{questao}

\begin{questao}{\nivel{N4} 07.48 — FGV}
Em processo de controle, o relatório aponta dano de R\$ 1,2 milhão, mas utiliza base de dados sem origem documentada e cujo histórico foi sobrescrito. Antes de imputar débito, a providência mais tecnicamente defensável é:
\begin{enumerate}[label=\Alph*)]
\item aceitar o valor porque foi calculado por servidor público.
\item fortalecer a evidência, reconstruindo a origem, integridade e metodologia da quantificação e assegurando contraditório sobre esses elementos.
\item aplicar multa em dobro para compensar a falta de evidência.
\item dispensar identificação de responsável.
\item converter automaticamente o processo em ação penal.
\end{enumerate}
\end{questao}

\begin{questao}{\nivel{N4} 07.49 — Cebraspe (C/E)}
A interpretação integrada da Constituição estadual, Lei Orgânica, Regimento Interno e instruções normativas é necessária porque competência, rito e obrigação operacional podem estar distribuídos em diplomas diferentes, respeitada a hierarquia normativa.
\end{questao}

\begin{questao}{\nivel{N4} 07.50 — FGV}
Ao revisar a conformidade de uma transferência por emenda parlamentar, o auditor encontra: (i) proponente identificado; (ii) plano de trabalho genérico; (iii) pagamento realizado; (iv) ausência de metas mensuráveis; (v) notas fiscais sem vínculo com etapa física; e (vi) portal que não permite localizar o processo administrativo. Qual conclusão é mais adequada?
\begin{enumerate}[label=\Alph*)]
\item A identificação do proponente é suficiente para afastar qualquer risco.
\item Há deficiências de planejamento, transparência e rastreabilidade que comprometem a demonstração de regular aplicação do recurso.
\item A existência de nota fiscal torna desnecessário verificar entrega.
\item A ausência de processo administrativo é irrelevante.
\item O Tribunal só poderia examinar a legalidade formal da emenda, nunca sua execução.
\end{enumerate}
\end{questao}
